\uuid{H7s3}
\titre{ Matrice dépendant d'un paramètre }

\niveau{L2}
\module{Algèbre}
\chapitre{Réduction d'endomorphisme}
\sousChapitre{Diagonalisation}
\cpge{mpsi-mp2i.determinants.d, pcsi.calcul-matriciel-systemes-lineaires.d, ptsi.calcul-matriciel-systemes-lineaires.d, bcpst1.matrices, mp-mpi.reduction-endomorphismes.c, pc.reduction.b, psi.reduction-endomorphismes.b, pt.reduction-endomorphismes.b, bcpst2.valeurs-propres-vecteurs-propres.a, bcpst2.valeurs-propres-vecteurs-propres.b}
\theme{}
\auteur{Quentin Liard}
\datecreate{ 2026-03-05}
\organisation{AMSCC}
\difficulte{2}



\contenu{

\texte{
On considère la matrice
\[
A_\alpha=
\begin{pmatrix}
\alpha & 1\\
0 & 2
\end{pmatrix}.
\]
}

\begin{enumerate}

\item \question{Calculer le polynôme caractéristique.}
\reponse{$(X-\alpha)(X-2)$.}

\item \question{Pour quelles valeurs de $\alpha$ la matrice est-elle diagonalisable ?}
\reponse{Pour $\alpha\neq2$ (valeurs propres distinctes).}

\item \question{Que se passe-t-il pour $\alpha=2$ ?}
\reponse{La matrice n’est pas diagonalisable car elle n’a qu’un seul vecteur propre.}

\end{enumerate}

}